\clearpage
\phantomsection
\addcontentsline{toc}{chapter}{Introduction}
\chapter*{INTRODUCTION}

\section*{1. Motivation}

The video game industry has grown into one of the largest entertainment sectors in the world, generating over \$180 billion in annual revenue \cite{yannakakis2018artificial}. Within this landscape, the \textit{Roguelike Deckbuilder} sub-genre has seen remarkable popularity, with titles such as \textit{Slay the Spire}, \textit{Monster Train}, and \textit{Inscryption} achieving both critical acclaim and millions of copies sold \cite{slaythespire}. These games are characterized by procedurally generated maps, permadeath mechanics, and an evolving card collection that grows organically through each run.

Despite their commercial success, Roguelike Deckbuilders are notoriously difficult to balance. A well-balanced instance must satisfy several competing desiderata simultaneously: the game should be challenging enough to remain engaging but not so punishing that players abandon it in frustration; the card pool should offer multiple viable build paths so that strategic variety is preserved; and the difficulty curve should scale smoothly across the game's fifteen or more floors so that no single floor constitutes an impassable wall. Achieving all of these properties simultaneously across the combinatorial space of possible card combinations, enemy encounters, relic synergies, and random deck draws is a task that defies purely manual tuning \cite{schell2019art}.

Traditional game-balancing practice relies on \textit{playtesting}: human players repeatedly attempt the game, report friction points, and designers manually adjust parameters. While essential, this process is slow, expensive, and bounded by the cognitive bandwidth of the design team. It is also inherently subjective: what feels balanced to an expert speedrunner may seem brutally unfair to a casual player. Automated simulation offers a route to complementing playtesting with data-driven optimization, but it raises its own challenge: what agent should be used to simulate the player?

This thesis addresses both the balancing problem and the agent problem together. We design and train a suite of \textbf{Reinforcement Learning (RL) agents} that learn to play the game and subsequently serve as the evaluators inside an \textbf{evolutionary optimization loop} that searches for well-balanced game configurations.

\section*{2. Problem Statement}

Let $\theta \in \mathbb{R}^n$ be a vector of game parameters (card damages, enemy health points, mana costs, room-type frequencies, etc.). Given a game simulator $G(\theta)$ and a population of player agents $\mathcal{A} = \{a_1, a_2, \ldots, a_k\}$, we wish to find $\theta^*$ such that the resulting game satisfies:
\begin{enumerate}
    \item \textbf{Balance}: The win rate $w(\theta, a)$ for a representative agent $a$ lies within a target range $[w_{\min}, w_{\max}]$, and the game does not trivialize or become unwinnable for any agent persona.
    \item \textbf{Engagement}: A large fraction of the available cards are ``viable'' -- i.e., choosing them correlates positively with winning -- supporting diverse build strategies.
    \item \textbf{Coherence}: No card is a dominant trap -- appearing attractive to the agent but consistently leading to defeat -- and parameter scalars stay within an interpretable range.
\end{enumerate}
Because these three objectives can conflict (making the game harder may increase balance but reduce engagement), we frame the search as a \textbf{multi-objective optimization problem} over the Pareto front of $(\text{Balance}, \text{Engagement}, \text{Coherence})$.

\section*{3. Research Objectives}

The concrete objectives of this thesis are:
\begin{enumerate}
    \item \textbf{Build a complete game prototype}: Implement a fully playable Roguelike Deckbuilder in C\# with turn-based combat, a procedurally generated map, a rich card and relic pool, and a data-driven parameter system loadable from JSON.
    \item \textbf{Establish a playability baseline}: Design and implement a rule-based \textit{Deterministic Heuristic Agent (DHA)} that confirms the game is playable and provides early balancing data.
    \item \textbf{Train four RL agent personas}: Replace the deterministic agent with four learned MaskablePPO agents representing distinct play styles (\textit{Aggressive}, \textit{Defensive}, \textit{Balanced}, \textit{Adaptive}), each trained with persona-specific reward shaping and Domain Randomization to generalize across varied game configurations.
    \item \textbf{Implement and compare two evolutionary optimizers}: Deploy a \textit{Single-Objective Genetic Algorithm} (GA) and the \textit{NSGA-II} multi-objective algorithm; compare their effectiveness at discovering balanced game configurations.
    \item \textbf{Analyze and visualize results}: Produce Pareto front analysis, card viability charts, and training curves to illuminate the behavior of both the optimization and the trained agents.
\end{enumerate}

\section*{4. Scope and Methodology}

The research scope covers the full software stack from game engine to optimization layer. The game prototype, while self-contained and playable, is designed as a research vehicle rather than a commercial product, and its parameter space is crafted to be representative of the real design challenge without requiring the scale of a production codebase.

Methodologically, the work proceeds through four phases:
\begin{enumerate}
    \item \textbf{Foundation}: Literature review on Roguelike games, evolutionary algorithms, reinforcement learning, and related automated balancing research.
    \item \textbf{System construction}: Implement the game engine (C\# .NET 9), the cross-platform integration layer (pythonnet), the Gymnasium RL environment, and the optimization modules.
    \item \textbf{Agent training}: Train the four MaskablePPO agents on Kaggle GPU infrastructure, iteratively refining reward functions and hyperparameters based on training logs.
    \item \textbf{Optimization and evaluation}: Run GA and NSGA-II experiments; collect, analyze, and visualize results.
\end{enumerate}

\section*{5. Main Contributions}

The primary contributions of this thesis are:
\begin{itemize}
    \item A \textbf{complete, open-source Roguelike Deckbuilder prototype} with a data-driven architecture that separates game logic from balance parameters, enabling fully automated parameter search.
    \item A \textbf{cross-platform integration bridge} between C\# game logic and Python RL training, enabling high-throughput simulation on Linux Kaggle GPU kernels without network overhead.
    \item \textbf{Four trained RL personas} for game evaluation, each exhibiting measurably distinct behaviors (win-rate range 24--34\%, statistically significant behavioral differences in HP expenditure and turn duration).
    \item An \textbf{empirical comparison of GA versus NSGA-II} for game balancing under identical simulation budgets, demonstrating that the hierarchical NSGA-II formulation converges 33\% faster while achieving superior scores on all three objectives.
    \item A \textbf{reusable framework design} applicable to other card-based or turn-based games with minor adaptation.
\end{itemize}

\section*{6. Thesis Structure}

The remainder of this thesis is organized as follows:

\begin{itemize}
    \item \textbf{Chapter 1 -- Background and Related Work}: Reviews the Roguelike Deckbuilder genre, game balancing challenges, Genetic Algorithms, NSGA-II, Reinforcement Learning theory (MDP, PPO, invalid action masking), Domain Randomization, and related prior work.
    \item \textbf{Chapter 2 -- System Analysis and Design}: Presents the three-tier architecture, the game engine design, the RL environment specification, the four agent persona designs, and the optimization layer design including genome representations and fitness functions.
    \item \textbf{Chapter 3 -- Implementation}: Details the C\# game engine implementation, cross-platform pythonnet integration, the Gymnasium environment wrapper, MaskablePPO training pipeline, and the GA/NSGA-II optimizer implementations.
    \item \textbf{Chapter 4 -- Experiments and Evaluation}: Reports experimental results for both the RL training phase and the evolutionary optimization phase, including ablation analysis and comparative discussion.
    \item \textbf{Chapter 5 -- Conclusion and Future Work}: Summarizes the achievements, discusses limitations, and outlines directions for future research.
\end{itemize}
